\begin{pastexam}[label={px:ncu-115-dsa-04}]{中央資工所 115 學年度 第 4 題}
（多選；每題 5 分，每個選項單獨計分，答錯每個選項倒扣 1 分）

Which of the following statements about greedy algorithms are true?
\begin{enumerate}[label=\Alph*.]
\item A greedy algorithm is correct only if the problem satisfies both the greedy-choice property and optimal substructure.
\item The fractional knapsack problem can be optimally solved using a greedy strategy, whereas the 0/1 knapsack problem cannot.
\item Dijkstra's shortest-path algorithm fails to produce correct results if the graph contains edges with negative weights.
\item Every problem that can be solved using dynamic programming can also be solved optimally using a greedy algorithm.
\item Prim's and Kruskal's algorithms are greedy algorithms for finding a minimum spanning tree.
\end{enumerate}
\end{pastexam}

\begin{pastsoln}
答案：A、B、E。

C 的「fails to produce correct results if the graph contains edges with negative weights」照字面是說
「只要有負邊就一定算錯」，本解照字面判斷為錯。 如果把它讀成「有負邊時不保證正確」，C 就是對的
（那是第 8 題 A 選項的說法）；作答時可以註明。

\textbf{A 對}。 資結書第~19.1.6~節：貪心法正確要靠兩個性質。 貪心選擇性質：第一步做的貪心選擇，
包含在某個最佳解裡；最佳子結構：做了這個選擇之後，剩下的子問題的最佳解和這個選擇合起來，就是
原問題的最佳解。 反過來說，貪心法對所有輸入都正確時，它的第一步一定落在某個最佳解裡，剩下的
各步又是在剩下的子問題上照同樣規則做，這兩個性質都成立。 證明貪心法正確，就是證明這兩件事。

\textbf{B 對}。 分數背包依「價值／重量」由大到小裝，最後一件裝一部分，得到最佳解（資結書定理~19.2）。
0/1 背包同樣的策略不對：容量 $50$，物品（重量, 價值）為 $(10, 60)$、$(20, 100)$、$(30, 120)$，
單位價值 $6, 5, 4$。 貪心先拿前兩件，總重 $30$，第三件放不下，價值 $160$；最佳是拿後兩件，總重
$50$，價值 $220$。 不論依哪一種標準排序，都有類似的反例，0/1 背包要用動態規劃。

\textbf{C 錯（照字面）}。 有負邊時 Dijkstra 演算法\textbf{可能}算錯（資結書第~25.2.5~節的例子），
但不是一定算錯。 例如只有一條邊 $s \to v$、權重 $-1$ 的圖，Dijkstra 演算法算出 $d[v] = -1$，
完全正確。 正確的說法是「有負邊時，Dijkstra 演算法不保證正確」。

\textbf{D 錯}。 0/1 背包可以用動態規劃求最佳解，貪心法不行（B 的反例）。 LCS、矩陣鏈乘也一樣。
能用動態規劃解的問題，多半只有最佳子結構，沒有貪心選擇性質。

\textbf{E 對}。 Kruskal 每一步選「不會形成環的最輕的邊」，Prim 每一步選「跨過目前的樹和其他頂點
之間最輕的邊」，都是只看眼前最好的選擇、選了不改，是貪心法；正確性來自切割性質（資結書第~24~章）。
\end{pastsoln}
